\documentclass[11pt]{article}
\usepackage[margin=1in]{geometry}
\usepackage[hidelinks]{hyperref}
\usepackage{amsmath,amssymb,mathrsfs}
\newcommand{\bishop}{\mathfrak{B}}

\title{Literature Status: Two Bishop-Property Questions}
\author{Run \texttt{fa\_banach\_001}, agent lane 13}
\date{2026-06-16}

\begin{document}
\maketitle

\section*{Status}

\textbf{literature\_already\_answered.}  Two open questions extracted from
arXiv:1211.2770 are answered in the later paper arXiv:1310.4035.  This
packet records provenance only; it is not an original proof from the run.

\section*{Source Questions}

The source paper is K. P. Hart, T. Kania, and T. Kochanek,
\emph{A chain condition for operators from \(C(K)\)-spaces},
arXiv:1211.2770.  In Section 3, Question 3.9 asks:
\[
\text{Does the class of compact Hausdorff spaces satisfying \((\bishop)\)
contain all Eberlein compact spaces?}
\]
It then asks more generally whether every compact Hausdorff space \(K\) for
which \(C(K)\) is Lindelof in its weak topology satisfies \((\bishop)\).

In Section 4, Question 4.3 asks whether the set of operators on a \(C(K)\)-space
which satisfy \((\bishop)\) is always a right ideal, and hence a two-sided ideal,
in \(\mathscr B(C(K))\).

\section*{Supporting Answers}

The supporting paper is T. Kania and R. J. Smith,
\emph{Chains of functions in \(C(K)\)-spaces}, arXiv:1310.4035.  Its abstract
states that it answers two questions arising after \((\bishop)\) was introduced.

First, arXiv:1310.4035 proves Theorem 1.1: if a class \(\mathscr D\) of compact
Hausdorff spaces is preserved under closed subspaces and Hausdorff quotients and
contains no non-metrizable linearly ordered space, then every member of
\(\mathscr D\) has \((\bishop)\).  The introduction then applies this theorem to
the class of compact spaces \(K\) for which \(C_p(K)\) is Lindelof and explicitly
says that this positively answers Question 3.9 from arXiv:1211.2770.  Since all
Eberlein compact spaces are Corson compact, and the paper recalls that Corson
compact spaces lie in that \(C_p(K)\)-Lindelof class, the Eberlein part is
answered affirmatively.

Second, the final section of arXiv:1310.4035 explicitly returns to Question 4.3
of arXiv:1211.2770 and gives a negative example.  With
\(B=B_{C(\beta\mathbb N)^*}\), it constructs a compact space
\(B\sqcup\beta\mathbb N\) and operators \(S,T\) such that \(T\) has
\((\bishop)\), while \(TS^{-1}=I_{C(B\sqcup\beta\mathbb N)}\) fails
\((\bishop)\).  Thus the Bishop operators need not form a right ideal.

\section*{Scope}

The Eberlein question is answered positively.  The more general weak-topology
wording from arXiv:1211.2770 is not identical to the \(C_p(K)\)-Lindelof class
handled in arXiv:1310.4035, but the supporting paper explicitly records the
positive answer to the Eberlein part and a broad pointwise-Lindelof strengthening.
The right-ideal question is answered negatively.

This packet does not settle Question 2.6 of arXiv:1211.2770 about set algebras
with the subsequential completeness property and vector measures satisfying the
rook condition.  A separate attempt note records that remaining hard core.

\section*{Search Evidence}

The run's cheap indexes were searched for arXiv:1211.2770, the paper title,
``Bishop'', ``right ideal'', ``Eberlein'', ``weakly Lindelof'', ``SCP'', and
``rook''.  No prior packet for arXiv:1211.2770 was found.  Web search for the
source title and the right-ideal/Eberlein phrases found arXiv:1310.4035, whose
abstract and introduction explicitly identify the source questions it answers.

\begin{thebibliography}{9}

\bibitem{HartKaniaKochanek2012}
K. P. Hart, T. Kania, and T. Kochanek,
\emph{A chain condition for operators from \(C(K)\)-spaces},
arXiv:1211.2770, 2012.
\url{https://arxiv.org/abs/1211.2770}.

\bibitem{KaniaSmith2013}
T. Kania and R. J. Smith,
\emph{Chains of functions in \(C(K)\)-spaces},
arXiv:1310.4035, 2013.
\url{https://arxiv.org/abs/1310.4035}.

\end{thebibliography}

\end{document}
